\chapter{Manage Store as a Service}

Catalog now starts and stops through its Blackship lifecycle without a
foreground console. We will apply the proven pattern to store. Its health check
requests the rendered product page, so one check crosses store, catalog, and
the database.

\section{Exercise 19A: Create the store launcher}

In the store console, stop the foreground server with \texttt{Ctrl-C}. If the
prompt belongs to the \texttt{store} user, return to \texttt{root}:

\begin{lstlisting}[language=sh]
exit
\end{lstlisting}

Create \texttt{/usr/local/sbin/store-server}:

\begin{lstlisting}[language=sh]
vi /usr/local/sbin/store-server
\end{lstlisting}

Enter:

\begin{lstlisting}[language=sh]
#!/bin/sh
set -eu

cd /usr/local/www/store
exec /usr/local/bin/ruby bin/rails server \
  --environment development \
  --binding 10.17.89.25 \
  --port 3000
\end{lstlisting}

Make the launcher executable and ask the shell to parse it without running it:

\begin{lstlisting}[language=sh]
chmod 0555 /usr/local/sbin/store-server
sh -n /usr/local/sbin/store-server
\end{lstlisting}

Expect no output from \texttt{sh -n}.

\section{Exercise 19B: Add the store service}

Create \texttt{/usr/local/etc/rc.d/store}:

\begin{lstlisting}[language=sh]
vi /usr/local/etc/rc.d/store
\end{lstlisting}

Enter:

\begin{lstlisting}[language=sh]
#!/bin/sh

# PROVIDE: store
# REQUIRE: NETWORKING
# KEYWORD: shutdown

. /etc/rc.subr

name="store"
desc="Rails store application"
rcvar="store_enable"

load_rc_config "$name"

: ${store_enable:="NO"}
: ${store_logfile:="/var/log/store.log"}

pidfile="/var/run/${name}.pid"
required_files="/usr/local/sbin/store-server"
command="/usr/sbin/daemon"
procname="/usr/sbin/daemon"
command_args="-f -r -P ${pidfile} -u store \
  -o ${store_logfile} \
  /usr/local/sbin/store-server"

run_rc_command "$1"
\end{lstlisting}

The \texttt{-f} option closes descriptors inherited from the Blackship hook
before \texttt{daemon} connects the service to its log. Without it, Puma can be
healthy while the hook waits indefinitely for its captured output pipes to
close.

Validate and start the service once:

\begin{lstlisting}[language=sh]
chmod 0555 /usr/local/etc/rc.d/store
sh -n /usr/local/etc/rc.d/store
service store onestart
service store onestatus
sockstat -4 -l | grep ':3000'
tail -n 20 /var/log/store.log
\end{lstlisting}

From the FreeBSD host, verify the page directly:

\begin{lstlisting}[language=sh]
fetch -qo - http://10.17.89.25:3000/ | \
  grep 'FreeBSD Handbook'
\end{lstlisting}

\section{Exercise 19C: Declare the store lifecycle}

Open the \texttt{store} jail's \textbf{Manifest} editor in BSDRunner. Add this
line with its top-level jail settings, before \texttt{[jails.network]}:

\begin{lstlisting}
depends_on = ["catalog"]
\end{lstlisting}

Append:

\begin{lstlisting}
[jails.healthcheck]
enabled = true

[[jails.healthcheck.checks]]
name = "store-page"
command = "fetch -qo /dev/null http://10.17.89.25:3000/"
target = "jail"
start_period = 30
interval = 30
timeout = 5
retries = 3

[[jails.hooks]]
phase = "post_start"
target = "jail"
command = "/usr/sbin/service"
args = ["store", "onestart"]
on_failure = "abort"

[[jails.hooks]]
phase = "pre_stop"
target = "jail"
command = "/usr/sbin/service"
args = ["store", "onestop"]
on_failure = "continue"
\end{lstlisting}

Run \textbf{Dry run}, review the dependency, health check, and hooks, then
choose \textbf{Apply}. Restart store through BSDRunner. If the Restart action
is unavailable, use Stop followed by Start; both paths execute the same
pre-stop and post-start hooks.

After store returns to the running state, verify its private entry point from
the FreeBSD host:

\begin{lstlisting}[language=sh]
fetch -qo - http://10.17.89.25:3000/ | \
  grep 'FreeBSD Handbook'
\end{lstlisting}

Then reload \texttt{http://HOST\_LAN\_IP:8080/} from another device on the
LAN. Use the host port selected when publishing the web edge if port 8080 was
already occupied. A request originating on the FreeBSD host does not prove the
external-interface PF redirect.

\section{Checkpoint}

Stop after the store page passes its Blackship health check and remains
available through nginx following the managed restart. Both Rails consoles can
now remain closed. The next exercise will turn the external request into a
repeatable black-box test.
